\documentclass{article} %210 mm × 297 mm
\usepackage[utf8]{inputenc}
\usepackage[a4paper,left=1.5cm, right=1.5cm, top=2cm, bottom=2cm]{geometry}
\usepackage{graphicx}
%\usepackage{blindtext}
\usepackage{multirow}
\usepackage{mhchem}
\usepackage{url}
\usepackage{subfigure}
\usepackage{amsfonts,latexsym} % para tener disponibilidad de diversos simbolos
\usepackage{enumerate}
%\usepackage{booktabs}
\usepackage{float}
%\usepackage{threeparttable}
\usepackage{array,colortbl}
%\usepackage{ifpdf}
%\usepackage{rotating}
\usepackage{cite}
\usepackage{stfloats}
\usepackage{url}
\usepackage{listings}
\usepackage{fancyhdr}

\usepackage{color}
%\usepackage{hyperref}
%\usepackage{wrapfig}
\usepackage{array}
%\usepackage{multirow}
%\usepackage{adjustbox}
%\usepackage{nccmath}
\usepackage[dvipsnames]{xcolor}


\newcommand{\titulo}{Abgabe 7; MfN II.}
\newcommand{\nombre}{Liliana Sanfilippo}
\newcommand{\FCE}{\textbf{SoSe 2025}}
\newcommand{\txc}[1]{\textcolor{blue}{#1}}
 

\usepackage{fancyhdr}
\pagestyle{fancy}
\fancyhead{}
\fancyfoot{}
\fancyhead[LO]{\small\nombre}
\fancyhead[RE]{\small\titulo}
\fancyhead[LO]{\small\FCE}
\fancyfoot[RO,LE] {\thepage}

\setcounter{page}{1} %Número de la página, la editorial lo cambiará



\usepackage[onehalfspacing]{setspace}
\begin{document}


\begin{tabular}{p{12cm}}
{{\Large\bf\titulo}}\\
\vspace{0.2cm}\\
 {\large\nombre}\\
\end{tabular} 
\vspace{0.2cm}\\
\section{Aufgabe 7.1}
Der Vektorraum $\mathbb{R}^4$ sei mit dem kanonischen inneren Produkt versehen \txc{(Es existiert ein Standardskalarprodukt)}. Bestimmen Sie eine Orthonormalbasis für den Untervektorraum $U$. 
\[
v_1 = \begin{pmatrix} 
1\\1\\0\\1
\end{pmatrix}, \quad v_2 = \begin{pmatrix} 
1\\-2\\0\\0
\end{pmatrix}, \quad v_3 = \begin{pmatrix} 
1\\0\\-1\\2
\end{pmatrix}
\]
\\ 
$v_1$ normalisieren: 
\[
w_1 = \frac{v_1}{||v_1||} = \frac{v_1}{\sqrt{(<v_1, v_1>)}}  = \frac{v_1}{\sqrt{1 + 1+ 0 + 1}} = \frac{v_1}{\sqrt{3}}   = \frac{1}{\sqrt{3}} \begin{pmatrix} 
1\\1\\0\\1
\end{pmatrix} =  \frac{\sqrt{3}}{3} \begin{pmatrix} 
1\\1\\0\\1
\end{pmatrix} = \begin{pmatrix} 
\frac{\sqrt{3}}{3}\\\frac{\sqrt{3}}{3}\\0\\\frac{\sqrt{3}}{3}
\end{pmatrix} 
\]
$v_2$ othogonalisieren 
\[
v_2' = v_2 - <w_1, v_2> \cdot  w_1 =  \begin{pmatrix} 
1\\-2\\0\\0
\end{pmatrix} - <\begin{pmatrix} 
\frac{\sqrt{3}}{3}\\\frac{\sqrt{3}}{3}\\0\\\frac{\sqrt{3}}{3}
\end{pmatrix} , \begin{pmatrix} 
1\\-2\\0\\0
\end{pmatrix}> \cdot \begin{pmatrix} 
\frac{\sqrt{3}}{3}\\\frac{\sqrt{3}}{3}\\0\\\frac{\sqrt{3}}{3}
\end{pmatrix}  = \begin{pmatrix} 
1\\-2\\0\\0
\end{pmatrix} - (\frac{\sqrt{3}}{3} - \frac{2 \sqrt{3}}{3}) \cdot \begin{pmatrix} 
\frac{\sqrt{3}}{3}\\\frac{\sqrt{3}}{3}\\0\\\frac{\sqrt{3}}{3}
\end{pmatrix} = 
\begin{pmatrix} 
1\\-2\\0\\0
\end{pmatrix} + \frac{\sqrt{3}}{3} \cdot \begin{pmatrix} 
\frac{\sqrt{3}}{3}\\\frac{\sqrt{3}}{3}\\0\\\frac{\sqrt{3}}{3}
\end{pmatrix} 
\]
\[
= 
\begin{pmatrix} 
1\\-2\\0\\0
\end{pmatrix} + \begin{pmatrix} 
\frac{1}{3}\\\frac{1}{3}\\0\\\frac{1}{3}
\end{pmatrix}  = \begin{pmatrix} 
\frac{4}{3}\\\frac{-5}{3}\\0\\\frac{1}{3}
\end{pmatrix} 
\]

$v_2'$ normalisieren 
\[
w_2 = \frac{v_2'}{||v_2'||} 
= v_2' \cdot \frac{1}{\sqrt{(\frac{4}{3})^2 + (\frac{-5}{3})^2 + (\frac{1}{3})^2}} 
= v_2' \cdot  \frac{1}{\sqrt{\frac{16}{9} + \frac{25}{9} + \frac{1}{9}}} 
= v_2' \cdot  \frac{1}{\sqrt{\frac{42}{9}}} 
= v_2' \cdot  \frac{1}{\frac{\sqrt{42}}{\sqrt{9}}} 
= v_2' \cdot  \frac{3}{\sqrt{42}} 
= v_2' \cdot  \frac{3\sqrt{42}}{42} 
= v_2' \cdot  \frac{\sqrt{42}}{14} 
\]
\[
= \begin{pmatrix} 
\frac{4}{3}\\\frac{-5}{3}\\0\\\frac{1}{3}
\end{pmatrix} \cdot \frac{\sqrt{42}}{14}  = 
\begin{pmatrix} 
\frac{4\sqrt{42}}{42}\\\frac{-5\sqrt{42}}{42}\\0\\\frac{\sqrt{42}}{42}
\end{pmatrix} = 
\begin{pmatrix} 
\frac{2\sqrt{42}}{21}\\\frac{-5\sqrt{42}}{42}\\0\\\frac{\sqrt{42}}{42}
\end{pmatrix}
\]
$v_3$ othogonalisieren 
\begin{align*}
    v_3' &= v_3 - <w_2, v_3> \cdot  w_2  -  <w_1, v_3> \cdot  w_1 \\
    & = \begin{pmatrix} 
1\\0\\-1\\2
\end{pmatrix} - <\begin{pmatrix} 
\frac{2\sqrt{42}}{21}\\\frac{-5\sqrt{42}}{42}\\0\\\frac{\sqrt{42}}{42}
\end{pmatrix}, \begin{pmatrix} 
1\\0\\-1\\2
\end{pmatrix}> \cdot\begin{pmatrix} 
\frac{2\sqrt{42}}{21}\\\frac{-5\sqrt{42}}{42}\\0\\\frac{\sqrt{42}}{42} 
\end{pmatrix}  -  <\begin{pmatrix} 
\frac{\sqrt{3}}{3}\\\frac{\sqrt{3}}{3}\\0\\\frac{\sqrt{3}}{3}
\end{pmatrix} , \begin{pmatrix} 
1\\0\\-1\\2
\end{pmatrix}> \cdot  \begin{pmatrix} 
\frac{\sqrt{3}}{3}\\\frac{\sqrt{3}}{3}\\0\\\frac{\sqrt{3}}{3}
\end{pmatrix}  \\
%
& = \begin{pmatrix} 
1\\0\\-1\\2
\end{pmatrix}  - (\frac{2\sqrt{42}}{21} \cdot 1 + 0 + 0 + 2 \cdot  \frac{\sqrt{42}}{42}) \cdot \begin{pmatrix} 
\frac{2\sqrt{42}}{21}\\\frac{-5\sqrt{42}}{42}\\0\\\frac{\sqrt{42}}{42}
\end{pmatrix} - (\frac{\sqrt{3}}{3} + 0 + 0 + \frac{2\sqrt{3}}{3}) \cdot  \begin{pmatrix} 
\frac{\sqrt{3}}{3}\\\frac{\sqrt{3}}{3}\\0\\\frac{\sqrt{3}}{3}
\end{pmatrix}  \\
&= \begin{pmatrix} 
1\\0\\-1\\2
\end{pmatrix} - \frac{3\sqrt{42}}{21} \cdot \begin{pmatrix} 
\frac{2\sqrt{42}}{21}\\\frac{-5\sqrt{42}}{42}\\0\\\frac{\sqrt{42}}{42}
\end{pmatrix} -  \sqrt{3} \cdot  \begin{pmatrix} 
\frac{\sqrt{3}}{3}\\\frac{\sqrt{3}}{3}\\0\\\frac{\sqrt{3}}{3}
\end{pmatrix}  \\
% 
&= \begin{pmatrix} 
1\\0\\-1\\2
\end{pmatrix} - \begin{pmatrix} 
\frac{4}{7}\\\frac{-5}{7}\\0\\\frac{1}{7}
\end{pmatrix} -  \begin{pmatrix} 
1 \\1\\0\\1
\end{pmatrix}   = \begin{pmatrix}
-\frac{4}{7} \\ - \frac{2}{7} \\ -1 \\ \frac{6}{7}
\end{pmatrix}
\end{align*}

$v_3'$ normalisieren. 
\[
w_3 = \frac{v_3'}{||v_3'||}  = v_3' \cdot \frac{1}{\sqrt{
(-\frac{4}{7})^2 + (-\frac{2}{7})^2 + (-1)^2 + (\frac{6}{7})^2}}  = v_3' \cdot \frac{1}{\sqrt{\frac{16}{49} + \frac{4}{49} + \frac{49}{49} + \frac{36}{49}}} = v_3' \cdot \frac{1}{\sqrt{\frac{105}{49}}} = v_3' \cdot \frac{\sqrt{105}}{15}  = \begin{pmatrix}
-\frac{4\sqrt{105}}{105} \\ -\frac{2\sqrt{105}}{105}  \\ -\frac{\sqrt{105}}{15}  \\ -\frac{2\sqrt{105}}{35} 
\end{pmatrix}
\]
Orthonormalbasis: 
\[
<
\begin{pmatrix} 
\frac{\sqrt{3}}{3}\\\frac{\sqrt{3}}{3}\\0\\\frac{\sqrt{3}}{3}
\end{pmatrix}, 
\begin{pmatrix} 
\frac{2\sqrt{42}}{21}\\\frac{-5\sqrt{42}}{42}\\0\\\frac{\sqrt{42}}{42}
\end{pmatrix},
\begin{pmatrix}
-\frac{4\sqrt{105}}{105} \\ -\frac{2\sqrt{105}}{105}  \\ -\frac{\sqrt{105}}{15}  \\ -\frac{2\sqrt{105}}{35} 
\end{pmatrix}
>
\]
\section{Aufgabe 7.2}
\subsection*{a)}
Zeigen Sie, dass die Matrix $B := A \overline{A}^T$ \txc{(also $B = AA^H$)} hermitesch ist; \txc{(zeigen, dass $B$ gleich der konjugiert transponierten Matrix $B^H$ ist)} \\
\[
B^H = (A A^H)^H = A^H (A^H)^H= A^H A = B
\]
\subsection*{b)}
Ist $A$ weiterhin auch invertierbar, dann ist $\beta_B (x, y) := x^T B \overline{y}$ \txc{(eine bilineare Form)} positiv definit; \txc{$\beta_B (x, y)$ ist positiv definit, wenn für alle $x \neq 0$ gilt $x^T A x > 0$}.\\
\[
\beta_B (x,x) = x^T A \overline{A}^T \overline{x} = x^T A \overline{A^T x} = (x A^T)^T \overline{A^T x} = \sum_{i=1}^n (A^T x)_i \overline{(A^T x)_i} = || A^T x ||^2
\]
Da $x \leq 0$ und negative Vorzeichen sich aufheben würden, muss $|| A^T x ||^2$ positiv sein.
\subsection*{c)}
/
\section{Aufgabe 7.3}
Auf $\mathbb{C}^3$ sei das innere Produkt $\gamma$ gegeben durch die Gram-Matrix
\[C := \begin{pmatrix}
2 &1 + i & -i \\
1 - i & 3 & 0 \\
i & 0& 4
\end{pmatrix}
\]
Bestimmen Sie eine bezüglich $\gamma$ orthogonale Basis des $\mathbb{C}^3$. \\ \\
Standardbasen orthogonalisieren mit $\gamma$:
\[
v_1 = e_1 = \begin{pmatrix}
1\\0\\0
\end{pmatrix},  e_2 = \begin{pmatrix}
0\\1\\0
\end{pmatrix},  e_3 = \begin{pmatrix}
0\\0\\1
\end{pmatrix}
\]
$e_2$ orthogonalisieren: 
\[
v_2 = e_2 - \frac{\gamma(e_2, v_1)}{\gamma(v_1, v_1)} v_1 = e_2 \frac{e_2^H C v_1}{v_1^H C v_1}v_1    = e_2 - \frac{C_{21}}{C_{11}}v_1  = e_2 - \frac{i-1}{2}v_1 
\]
\[
= \begin{pmatrix}
0\\1\\0
\end{pmatrix} -  \frac{i-1}{2} \cdot \begin{pmatrix}
1\\0\\0
\end{pmatrix} = \begin{pmatrix}
0\\1\\0
\end{pmatrix} -  \begin{pmatrix}
\frac{i-1}{2}\\0\\0
\end{pmatrix} = \begin{pmatrix}
-\frac{i-1}{2}\\1\\0
\end{pmatrix} 
\]
$e_3$ orthogonalisieren: 

\[
v_2^H C v_2 = \left(\begin{matrix}
\frac{1-i}{2} & 1 & 0
\end{matrix}\right) C \begin{pmatrix}
-\frac{i-1}{2}\\1\\0
\end{pmatrix}  = \left(\begin{matrix}
2-2*i & 4 & \frac{-1-i}{2}
\end{matrix}\right) \begin{pmatrix}
-\frac{i-1}{2}\\1\\0
\end{pmatrix} = 4-2i
\]
\[
e_3^H C v_2 = \left(\begin{matrix}
0&0&1
\end{matrix}\right) C \begin{pmatrix}
-\frac{i-1}{2}\\1\\0
\end{pmatrix}  = \left(\begin{matrix}
i & 0 & 4
\end{matrix}\right) \begin{pmatrix}
-\frac{i-1}{2}\\1\\0
\end{pmatrix} = \frac{1+i}{2}
\]
\[
v_3 = e_3 - \frac{\gamma(e_3, v_1)}{\gamma(v_1, v_1)} v_1 - \frac{\gamma(e_3, v_2)}{\gamma(v_2, v_2)} v_2 = e_3 - \frac{C_{31}}{C_{11}} v_1 - \frac{e_3^H C v_2}{v_2^H C v_2}v_2  = e_3 - \frac{i}{2} v_1 - \frac{\frac{1+i}{2}}{4-2i} v_2 = e_3 - \frac{i}{2} v_1 - \frac{{1+i}}{8-4i} v_2
%= e_3 - \frac{i}{2} v_1 - \frac{e_3^H C e_2 - e_3^H C \frac{1-i}{2} v_1}{v_2^H C v_2}
\]
\[
 = \begin{pmatrix}
0\\0\\1
\end{pmatrix} - \frac{i}{2} \begin{pmatrix}
1\\0\\0
\end{pmatrix} - \frac{{1+i}}{8-4i}  \begin{pmatrix}
-\frac{i-1}{2}\\1\\0
\end{pmatrix} = \begin{pmatrix}
- \frac{i}{2}\\0\\1
\end{pmatrix} 
+ \begin{pmatrix}
\frac{i^2-1}{16-8i}\\- \frac{{1+i}}{8-4i} \\0
\end{pmatrix} = \left(\begin{matrix}
\frac{-2-11i}{20} \\
\frac{-1-3i}{20} \\
1
\end{matrix}\right)
\]
Die Basis ist; 

\[
 \begin{pmatrix}
1\\0\\0
\end{pmatrix}, 
\begin{pmatrix}
-\frac{i-1}{2}\\1\\0
\end{pmatrix}  , \left(\begin{matrix}
\frac{-2-11i}{20} \\
\frac{-1-3i}{20} \\
1
\end{matrix}\right)
\]
\section{Aufgabe 7.4}
/
\end{document}